

\section{Appendix}

In this appendix, we provide additional analysis and visualizations of the debates used in the main paper in \sect{sect:additional_results}. We further provide detailed experimental details on each dataset in \sect{sect:evaluation}.

\subsection{Additional Results}
\label{sect:additional_results}


\begin{wrapfigure}{r}{6.0cm}
    \vspace{-20pt}
    \centering
    \includegraphics[width=\linewidth]{fig/consensus_agents.pdf}
    \vspace{-15pt}
    \caption{\textbf{Effect of Prompts on Consensus.} Using a short debate prompt induces faster consensus between agents}
    \vspace{-15pt}
    \label{fig:consensus}
\end{wrapfigure}

\vspace{3pt}
\myparagraph{Consensus Between Agents.} In \fig{fig:consensus}, we illustrate the consensus between agents using either short or long consensus prompts discussed in \fig{tbl:prompt}. The use of debate prompts that encourage agents to adapt more to the opinions of other agents improves consensus. 



\myparagraph{Additional Qualitative Visualizations.} We added additional qualitative visualizations of the debate process. In \fig{fig:gsm-1}, \fig{fig:gsm-2}, \fig{fig:gsm-3}, \fig{fig:gsm-4}, \fig{fig:gsm-5}, we illustrate debates between agents in the GSM8K dataset which result in the correct answer. In \fig{fig:gsm-6}, \fig{fig:gsm-7}, \fig{fig:gsm-8}, we further illustrate debates in GSM8K which lead to the incorrect answer. We further provide an example illustration of debate in arithmetic in \fig{fig:result_math}, arithmetic with summarization of individual responses of agents in \fig{fig:result_math_summarize}, MMLU in \fig{fig:result_mmlu}, a debate with the full contents biographies in \fig{fig:result_biography}, and debate in chess in \fig{fig:result_chess}. In general, we found that debate improved the performance of final generated answers, though sometimes answers would converge to the incorrect value.

\subsection{Evaluation Details}
\label{sect:evaluation}

We provided detailed evaluation details for each setting in the paper. We run all experiments using the \texttt{gpt-3.5-turbo-0301} model. We provide a table listing the prompts used to prompt models and initialize debate in \tbl{tbl:prompt_appendix}.

\myparagraph{Arithmetic.} To evaluate the arithmetic task, we generated six random integers for each task between 0 and 30. We then evaluated the extent to which the correct integer answer was correctly obtained. We evaluated models on one hundred generated arithmetic tasks.

\myparagraph{Grade School Math.} To evaluate the GSM8K task, we evaluated the accuracy at which models were able to obtain the final correct answer, as extracted from a box. We evaluated models on one hundred grade school math problems.

\myparagraph{Chess.} To evaluate the chess reasoning task, we used chess games from \url{https://www.pgnmentor.com/players/Adams.zip}. We asked chatGPT to predict the next move for white to move at turn 14 and  reported the relative Stockfish pawn score with search depth 20 after executing the suggested move from chatGPT.  We evaluated models on three hundred selected chess games.

\myparagraph{Biographies.} To evaluate the biographies task, we compare each generated bullet point biography for a person with a ground truth set of facts about the person extracted from Wikipedia. We iteratively loop through each ground truth fact, and validate the extent to which the generated biography matches a particular bullet by prompting chatGPT with the prompt: \emph{Consider the following biography of <person>: <generated biography>  Is the above biography above consistent with the fact below? <ground truth bullet> Give a single-word answer, yes, no, or uncertain.} We then evaluate and report the percentage of ground bullets that chatGPT returns either yes or no on. We ignored ground truth bullets that chatGPT returns returned uncertain.

We found this evaluation metric provided a fast way to evaluate how relatively correct a generated bullet point biography is. However, we found that generated facts could contain incorrect information that was not captured in the ground truth bullet and thus could not be validated through this metric. Nevertheless, we believe this evaluation scheme estimates the relative accuracy of a generated biography.

\myparagraph{MMLU.} To evaluate MMLU, we measured the accuracy in which models were able to select the correct multiple-choice answer in each problem. We evaluated models on one hundred selected MMLU questions randomly distributed across each of the subject areas.

\myparagraph{Chess Validity.} To evaluate chess validity, we consider the BIG-Bench Chess-State Tracking Benchmark ~\cite{srivastava2022beyond}, where we used the hardest reported task in the benchmark \texttt{synthetic\_short}. Each generated answer was deemed correct as long as it was one of the valid answers in the sequence. We evaluated models of one hundred selected chess validity tasks.

\begin{table*}[t]
\small\setlength{\tabcolsep}{5.5pt}
\centering
\scalebox{0.8}{

\begin{tabular}{c|c|l}
     {\bf Task} & Type & {\bf Prompt} \\
      \midrule
     \multirow{4}{*}{Arithmetic}  & Starting & \emph{What is the result of \{\}+\{\}*\{\}+\{\}-\{\}*\{\}? Make sure to state your answer at the end of the response.} \\
     \cmidrule{2-3}
      & \multirow{3}{*}{Debate} & \emph{These are the recent/updated opinions from other agents: <other agent responses>  Use these opinions} \\
      & & \emph{carefully as additional advice, can you provide an updated answer? Make sure to state your answer} \\
      & & \emph{at the end of the response.} \\
      \midrule
      \multirow{6}{*}{GSM8K}  & \multirow{2}{*}{Starting} & \emph{Can you solve the following math problem? <Problem> Explain your reasoning. Your final answer }  \\
      & & \emph{should be a single numerical number, in the form \textbackslash boxed\{\{answer\}\}, at the end of your response.} \\
      \cmidrule{2-3}
     & \multirow{4}{*}{Debate} & \emph{These are the solutions to the problem from other agents: <other agent responses> Using the solutions} \\
     & & \emph{from other agents as additional information, can you provide your answer to the math problem? The original } \\
     & & \emph{math problem is <Problem>. Your final answer should be a single numerical number, in the form} \\
     & & \emph{ \textbackslash boxed\{\{answer\}\}, at the end of your response.} \\
      \midrule
      \multirow{2}{*}{Chess} & \multirow{3}{*}{Starting} & \emph{Here is the current sequence of moves in a chess game: <moves>. What is the best chess move I should  } \\
       &  & \emph{ execute next? Give a single move suggestion of the form 14. <XXX> and make sure the chess move } \\
       & & \emph{is valid in the current board state.} \\
     \cmidrule{2-3}
       & \multirow{3}{*}{Debate} & \emph{Here are other chess move suggestions from other agents: <other agent responses> Using the chess suggestions} \\
       & & \emph{ from other agents as additional advice and your earlier generated solution,  can you give me your updated thoughts  } \\
       & & \emph{ on the best next chess move I should play given the chess sequence {}? Give a single move suggestion of the form} \\
       & & \emph{ 14. <XXX> and make sure the chess move is valid in the current board state.} \\
      \midrule
      \multirow{4}{*}{Biographies} & \multirow{2}{*}{Starting} & \emph{Give a bullet point biography of {} highlighting their contributions and achievements as a computer scientist,} \\
      & & \emph{with each fact separated with a new line character.} \\
     \cmidrule{2-3}
       & \multirow{2}{*}{Debate} & \emph{Here are some bullet point biographies of <person> given by other agents: <other agent response> Closely} \\
       & & \emph{ examine your biography and the biography of other agents and provide an updated bullet point biography.} \\
      \midrule
      \multirow{5}{*}{MMLU} & \multirow{2}{*}{Starting} & \emph{Can you answer the following question as accurately as possible? {}: A) {}, B) {}, C) {}, D) {} Explain your answer, } \\
      & & \emph{putting the answer in the form (X) at the end of your response.} \\
     \cmidrule{2-3}
       & \multirow{3}{*}{Debate} & \emph{These are the solutions to the problem from other agents: <other agent responses> Using the reasoning}\\
       & & \emph{ from other agents as additional advice, can you give an updated answer? Examine your solution and}\\
       & & \emph{ that other agents. Put your answer in the form (X) at the end of your response.} \\
      \midrule
      \multirow{6}{*}{Chess Validity} & \multirow{3}{*}{Starting} & \emph{Given the chess game {}, give one valid destination square for the chess piece at {}. State the destination square } \\
      & & \emph{in the form (X), where X follows the regex [a-h][1-8], for example (e5). Give a one line explanation } \\
      & & \emph{of why your destination square is a valid move.} \\
     \cmidrule{2-3}
       & \multirow{3}{*}{Debate} & \emph{Here are destination square suggestions from other agents: <other agent responses>  Can you double}\\
       & & \emph{check that your destination square is a valid move? Check the valid move justifications from other agents.}\\
       & & \emph{State your final answer in a newline with a 2 letter response following the regex [a-h][1-8].} \\
    \bottomrule
\end{tabular}
}
\captionof{figure}{\small {\bf Prompts in each task.} List of prompts used in each task }
\label{tbl:prompt_appendix}
\vspace{-10pt}
\end{table*}


\begin{figure}[t]
    \centering
    \includegraphics[width=\linewidth]{fig/gsm-1.pdf}
    \caption{\textbf{Example of a correct GSM8K Debate.}}
    \label{fig:gsm-1}
    \vspace{-10pt}
\end{figure}

\begin{figure}[t]
    \centering
    \includegraphics[width=\linewidth]{fig/gsm-2.pdf}
    \caption{\textbf{Example of Correct GSM8K Debate.}}
    \label{fig:gsm-2}
    \vspace{-10pt}
\end{figure}

\begin{figure}[t]
    \centering
    \includegraphics[width=\linewidth]{fig/gsm-3.pdf}
    \caption{\textbf{Example of Correct GSM8K Debate.}}
    \label{fig:gsm-3}
    \vspace{-10pt}
\end{figure}

\begin{figure}[t]
    \centering
    \includegraphics[width=\linewidth]{fig/gsm-4.pdf}
    \caption{\textbf{Example of Correct GSM8K Debate.}}
    \label{fig:gsm-4}
    \vspace{-10pt}
\end{figure}

\begin{figure}[t]
    \centering
    \includegraphics[width=\linewidth]{fig/gsm-5.pdf}
    \caption{\textbf{Example of Correct GSM8K Debate.}}
    \label{fig:gsm-5}
    \vspace{-10pt}
\end{figure}


\begin{figure}[t]
    \centering
    \includegraphics[width=\linewidth]{fig/gsm-6.pdf}
    \caption{\textbf{Example of Incorrect GSM8K Debate.}}
    \label{fig:gsm-6}
    \vspace{-10pt}
\end{figure}

\begin{figure}[t]
    \centering
    \includegraphics[width=0.95\linewidth]{fig/gsm-7.pdf}
    \caption{\textbf{Example of Incorrect GSM8K Debate.}}
    \label{fig:gsm-7}
    \vspace{-10pt}
\end{figure}

\begin{figure}[t]
    \centering
    \includegraphics[width=\linewidth]{fig/gsm-8.pdf}
    \caption{\textbf{Example of Incorrect GSM8K Debate.}}
    \label{fig:gsm-8}
    \vspace{-10pt}
\end{figure}


\begin{figure}[t]
    \centering
    \includegraphics[width=\linewidth]{fig/math.pdf}
    \caption{\textbf{Example of Arithmetic Debate.}}
    \label{fig:result_math}
    \vspace{-10pt}
\end{figure}

\begin{figure}[t]
    \centering
    \includegraphics[width=\linewidth]{fig/math.pdf}
    \caption{\textbf{Example of Arithmetic Debate with Summarization.} Four separate agents participate in debate, with two illustrated above. Instruction contains the summarized responses across agents.}
    \label{fig:result_math_summarize}
    \vspace{-10pt}
\end{figure}

\begin{figure}[t]
    \centering
    \includegraphics[width=\linewidth]{fig/mmlu.pdf}
    \caption{\textbf{Example of MMLU Debate.}}
    \label{fig:result_mmlu}
    \vspace{-10pt}
\end{figure}

\begin{figure}[t]
    \centering
    \includegraphics[width=\linewidth]{fig/biography.pdf}
    \caption{\textbf{Example of Biography Debate.} While we found that generated biographies after debate to be more accurate, many facts remain incorrect.}
    \label{fig:result_biography}
    \vspace{-10pt}
\end{figure}

\begin{figure}[t]
    \centering
    \includegraphics[width=\linewidth]{fig/result_chess.pdf}
    \caption{\textbf{Example of Chess Debate.}}
    \label{fig:result_chess}
    \vspace{-10pt}
\end{figure}




